\section{The article's eleven-vertex seed and its extension}\label{sec:article}\label{sec:example}\label{sec:seed3k8}
The seed $F_1=A$ is given in Table~\ref{ker:tab:seed}. Here we record the
finite reductions and Euler obstructions that establish the strict levels
of the $3k+8$ family.
Its profiles and type graph are
\begin{equation}\label{eq:seed-profiles}
 P_0=V(A)\setminus\{2\},\qquad P_1=V(A)\setminus\{1\},\qquad R=K_2.
\end{equation}
The first extension $F_2=B$ adds $p,x,y$, with $p$ adjacent to every seed
vertex and to neither auxiliary. The vertex $x$ is adjacent to $P_0$, the
vertex $y$ to $P_1$, and $xy$ is an edge. Thus $A$ has eleven vertices and
thirty-three edges, while $B$ has fourteen vertices and sixty-five edges.
Figure~\ref{fig:extension} shows the added edges in this first extension.

\begin{figure}[htbp]
\centering\begin{tikzpicture}[x=1cm,y=1cm]
\coordinate (w) at (0,0);
\coordinate (p) at (0,2);
\coordinate (x) at (-3,.6);
\coordinate (y) at (3,.6);
\coordinate (a) at (-3,-1.6);
\coordinate (b) at (3,-1.6);
\draw[jfef new edge] (p)--(w);
\draw[jfef new edge] (x)--(w);
\draw[jfef new edge] (y)--(w);
\draw[jfef new edge] (x)--(a);
\draw[jfef new edge] (y)--(b);
\draw[jfef new edge] (x) .. controls (-1.7,-1.7) and (1.7,-1.7) .. (y);
\draw[jfef new edge] (p) .. controls (-5,2.2) and (-5,-1.2) .. (a);
\draw[jfef new edge] (p) .. controls (5,2.2) and (5,-1.2) .. (b);
\node[ellipse,draw=black!70,fill=white,minimum width=3cm,minimum height=1.3cm,jfef label font] at (w) {$W$};
\node[jfef centre] at (p) {$p$};
\node[jfef auxiliary] at (x) {$x$};
\node[jfef auxiliary] at (y) {$y$};
\node[jfef vertex] at (a) {$1$};
\node[jfef vertex] at (b) {$2$};
\end{tikzpicture}

\caption{The thirty-two added edges. The ellipse $W$ represents
$V(A)\setminus\{1,2\}$, the set $Z$ of Section~\ref{sec:kernels}; an edge to $W$ means adjacency to all nine vertices.
Edges within $A$, including $\{1,2\}$, are specified by Table~\ref{ker:tab:seed}.}\label{fig:extension}
\end{figure}

\Needspace{6\baselineskip}
\subsection{The seed is strictly \texorpdfstring{$1$}{1}-dismantlable}
We will use folds in two different graphs. To show that \(A\) belongs
to level one, we fold each neighbourhood graph needed in a reduction.
To exclude level zero, we examine the folds available in \(A\) itself.

\begin{lemma}\label{lem:seed-in-D1}
The graph \(A\) belongs to \(\D_1\).
\end{lemma}
\begin{proof}
We delete the vertices of \(A\) in the order
\[
 5,2,1,4,0,6,7,8,9,3,
\]
leaving \(10\). Each row of Table~\ref{tab:seed-reduction} justifies
one deletion. The second column lists the neighbours of the vertex to
be deleted, after the preceding deletions from \(A\). The third column
gives the folds in the graph induced by those neighbours, and the last
column gives the surviving vertex.

To check a listed fold \(u\to t\), let \(H\) be the neighbourhood
graph for that row after its preceding folds. The adjacencies in
Table~\ref{tab:seed-adjacency} give \(N_H[u]\subseteq N_H[t]\), so the
fold is valid by Lemma~\ref{lem:dom}. In the final two rows the
neighbourhood graph is already a single vertex. Thus each row gives a
reduction of its neighbourhood graph to one vertex, which puts that
graph in \(\D_0\) by Definition~\ref{def:hierarchy}. The ten deletions
therefore dismantle \(A\) at level one.
\end{proof}

\begin{table}[htbp]
\centering\small\setlength{\tabcolsep}{4pt}
\caption{A complete level-one reduction of \(A\).}\label{tab:seed-reduction}
\begin{tabular}{cclc}
\toprule
Deleted vertex&Current neighbourhood&Folds within that neighbourhood&Survivor\\\midrule
$5$&$\{1,2,6,8,9\}$&$1\to2,\ 6\to2,\ 8\to2,\ 2\to9$&$9$\\
$2$&$\{1,6,7,8,9,10\}$&$6\to1,\ 7\to1,\ 8\to10,\ 9\to1,\ 1\to10$&$10$\\
$1$&$\{4,6,7,9,10\}$&$6\to10,\ 7\to4,\ 9\to4,\ 4\to10$&$10$\\
$4$&$\{0,7,8,9,10\}$&$9\to7,\ 7\to0,\ 0\to8,\ 8\to10$&$10$\\
$0$&$\{3,6,7,8\}$&$6\to3,\ 7\to3,\ 3\to8$&$8$\\
$6$&$\{3,10\}$&$3\to10$&$10$\\
$7$&$\{3,9\}$&$3\to9$&$9$\\
$8$&$\{3,10\}$&$3\to10$&$10$\\
$9$&$\{3\}$&\text{none}&$3$\\
$3$&$\{10\}$&\text{none}&$10$\\
\bottomrule
\end{tabular}

\end{table}
\FloatBarrier

\Needspace{7\baselineskip}
\begin{lemma}\label{lem:seed-not-D0}
The graph \(A\) does not belong to \(\D_0\).
\end{lemma}
\begin{proof}
To determine whether a level-zero reduction can start, we list the
folds available in \(A\). By Lemma~\ref{lem:dom}, a fold of \(u\) onto
\(t\) in the current graph \(H\) is permitted exactly when \(u\ne t\)
and \(N_H[u]\subseteq N_H[t]\). Comparing the closed neighbourhoods from
Table~\ref{tab:seed-adjacency} gives the complete lists below, first
in \(A\) and then after each forced deletion:
\[
\begin{array}{c|c}
 H&\text{all permitted folds}\\\hline
 A&5\to2\\
 A-5&9\to7\\
 A-\{5,9\}&\text{none}.
\end{array}
\]
Any level-zero reduction would have to delete \(5\) and then \(9\).
The remaining graph on \(\{0,1,2,3,4,6,7,8,10\}\) admits no fold,
so the reduction stops before reaching \(K_1\).
\end{proof}

\subsection{A dismantlable remainder}
The neighbourhood graph of \(p\) is \(A\), so \(p\) can be deleted at
level two. Theorem~\ref{thm:template} uses the hypothesis that the
remaining graph \(B-p\) can be dismantled at level one. In fact,
ordinary folds already suffice.

\begin{lemma}\label{lem:seed-rest}
The graph \(B-p\) belongs to \(\D_0\).
\end{lemma}
\begin{proof}
After removing \(p\), the remaining vertices are \(0,\ldots,10,x,y\).
Perform the following twelve folds in order:
\begin{equation}\label{eq:seed-rest-folds}
\begin{aligned}
 &0\to x,\quad3\to x,\quad4\to x,\quad7\to9,\quad8\to y,\quad9\to5,\\
 &1\to6,\quad2\to6,\quad5\to6,\quad6\to10,\quad10\to x,\quad x\to y.
\end{aligned}
\end{equation}
At each step, the adjacency table and the specified adjacencies
of \(x,y\) give the closed-neighbourhood inclusion for the indicated
fold in the graph left by the preceding folds. Only \(y\) survives.
Thus \(B-p\in\D_0\subseteq\D_1\), giving the required hypothesis of
Theorem~\ref{thm:template} for \(s=1\).
\end{proof}

\Needspace{5\baselineskip}
\subsection{Euler obstructions in the first extension}
A first deletion at level one would require a neighbourhood graph
in \(\D_0\). The centre fails this test because its neighbourhood
graph is \(A\). For each other vertex, we use Lemma~\ref{lem:qobstruction}:
a non-zero \(q\) excludes its neighbourhood graph from \(\D_0\).

\begin{lemma}\label{lem:seed-euler}
Every non-central vertex \(v\) of \(B\) satisfies
\(q(B[N_B(v)])\ne0\).
\end{lemma}
\begin{proof}
For this clique table only, write \(H_v=B[N_B(v)]\) for each of the thirteen vertices \(v\ne p\). This local notation follows the article table and is unrelated to the family \(H_k\) of Section~\ref{sec:family4}.
To evaluate \(q(H_v)\), we count the cliques of \(H_v\) by size.
In Table~\ref{tab:seed-cliques}, \(c_r(H_v)\) counts the cliques of
size \(r\). The empty clique gives \(c_0(H_v)=1\), and none of the
neighbourhood graphs has a clique of size greater than four.

To obtain a row directly from the adjacency data, enumerate the
subsets of \(N_B(v)\) and count those in which every pair of vertices
is adjacent. These counts are the coefficients of \(C_{H_v}\).
Alternatively, compute the same polynomial using the recurrence in
Lemma~\ref{lem:qrec}, starting from \(C_{\varnothing}(z)=1\).
Each step uses graphs with fewer vertices, so the recursion terminates.
In either case, evaluation at \(z=-1\) gives the last column of
Table~\ref{tab:seed-cliques}, where every entry is non-zero.
\end{proof}

\begin{table}[htbp]
\centering
\caption{All clique counts of \(H_v=B[N_B(v)]\) for \(v\ne p\).}\label{tab:seed-cliques}
\begin{tabular}{crrrrrr}
\toprule
\(v\)&\(\kappa_0(B_v)\)&\(\kappa_1(B_v)\)&\(\kappa_2(B_v)\)&\(\kappa_3(B_v)\)&\(\kappa_4(B_v)\)&\(q(B_v)\)\\\midrule
$0$&1&8&21&20&5&$-1$\\
$1$&1&9&25&24&6&$-1$\\
$2$&1&9&24&21&4&$-1$\\
$3$&1&9&25&24&6&$-1$\\
$4$&1&9&25&24&6&$-1$\\
$5$&1&8&20&16&2&$-1$\\
$6$&1&9&24&21&4&$-1$\\
$7$&1&9&25&25&7&$-1$\\
$8$&1&9&24&21&4&$-1$\\
$9$&1&9&25&24&6&$-1$\\
$10$&1&9&25&23&4&$-2$\\
$x$&1&11&35&37&11&$-1$\\
$y$&1&11&35&36&10&$-1$\\
\bottomrule
\end{tabular}

\end{table}

\subsection{The family at every strict level}
The checks on \(A\) and \(B\) now let us apply Theorem~\ref{thm:template}
to any number of layers, including zero layers.

\begin{theorem}\label{thm:family3k8}
For every integer \(k\geq1\), the graph \(F_k=G_{k-1}\)
with \(d=2\) and \(R=K_2\), specified by Table~\ref{tab:seed-adjacency},
\eqref{eq:seed-profiles}, and the adjacency rules of
Section~\ref{sec:construction} satisfies
\[
 F_k\in\D_k\setminus\D_{k-1},\qquad |V(F_k)|=3k+8.
\]
\end{theorem}
\begin{proof}
Lemmas~\ref{lem:seed-in-D1} and~\ref{lem:seed-not-D0} establish
\(A\in\D_1\setminus\D_0\). The second hypothesis follows from
\(B-p\in\D_0\subseteq\D_1\) by Lemma~\ref{lem:seed-rest}, and the
last is exactly Lemma~\ref{lem:seed-euler}.

Apply Theorem~\ref{thm:template} with \(s=1\) and \(d=2\).
It gives strict level \(m+1\) and order \(11+3m\) after \(m\) layers,
for every \(m\geq0\). Taking \(m=k-1\) gives level \(k\) and order
\(11+3(k-1)=3k+8\). The first two graphs are \(F_1=A\) and \(F_2=B\).
\end{proof}

\runcheck{article}
The report checks all ten outer deletions and twenty-three inner folds of the
seed, the forced folds excluding level zero, the twelve folds of $B-p$, and
all thirteen clique rows. It also checks the exact incidences drawn in the
figures: the eight schematic incidences in Figure~\ref{fig:extension} expand to thirty-two added edges. Figure~\ref{fig:complement} displays the seed by its twenty-two missing edges; Figure~\ref{fig:layers} shows all thirty edges on the nine vertices added in three layers.

\begin{figure}[htbp]
\centering\begin{tikzpicture}[x=1cm,y=1cm]
\path[use as bounding box] (-3.35,-2.9) rectangle (3.35,2.9);
\coordinate (v0) at (0.000000,2.350000);
\coordinate (v1) at (1.270506,1.976946);
\coordinate (v2) at (2.137635,0.976225);
\coordinate (v3) at (2.326080,-0.334440);
\coordinate (v4) at (1.776011,-1.538923);
\coordinate (v5) at (0.662072,-2.254808);
\coordinate (v6) at (-0.662072,-2.254808);
\coordinate (v7) at (-1.776011,-1.538923);
\coordinate (v8) at (-2.326080,-0.334440);
\coordinate (v9) at (-2.137635,0.976225);
\coordinate (v10) at (-1.270506,1.976946);
\draw[jfef complement edge] (v0) -- (v1);
\draw[jfef complement edge] (v0) -- (v2);
\draw[jfef complement edge] (v0) -- (v5);
\draw[jfef complement edge] (v0) -- (v9);
\draw[jfef complement edge] (v0) -- (v10);
\draw[jfef complement edge] (v1) -- (v3);
\draw[jfef complement edge] (v1) -- (v8);
\draw[jfef complement edge] (v2) -- (v3);
\draw[jfef complement edge] (v2) -- (v4);
\draw[jfef complement edge] (v3) -- (v4);
\draw[jfef complement edge] (v3) -- (v5);
\draw[jfef complement edge] (v4) -- (v5);
\draw[jfef complement edge] (v4) -- (v6);
\draw[jfef complement edge] (v5) -- (v7);
\draw[jfef complement edge] (v5) -- (v10);
\draw[jfef complement edge] (v6) -- (v7);
\draw[jfef complement edge] (v6) -- (v8);
\draw[jfef complement edge] (v6) -- (v9);
\draw[jfef complement edge] (v7) -- (v8);
\draw[jfef complement edge] (v7) -- (v10);
\draw[jfef complement edge] (v8) -- (v9);
\draw[jfef complement edge] (v9) -- (v10);
\node[jfef vertex] at (v0) {$0$};
\node[jfef vertex] at (v1) {$1$};
\node[jfef vertex] at (v2) {$2$};
\node[jfef vertex] at (v3) {$3$};
\node[jfef vertex] at (v4) {$4$};
\node[jfef vertex] at (v5) {$5$};
\node[jfef vertex] at (v6) {$6$};
\node[jfef vertex] at (v7) {$7$};
\node[jfef vertex] at (v8) {$8$};
\node[jfef vertex] at (v9) {$9$};
\node[jfef vertex] at (v10) {$10$};
\end{tikzpicture}

\caption{The complement of the eleven-vertex seed, with the twenty-two
non-edges listed by the adjacency data.}\label{fig:complement}
\end{figure}
\begin{figure}[htbp]
\centering\begin{tikzpicture}[x=1cm,y=1cm]
\path[use as bounding box] (-3.1,-2.7) rectangle (3.1,2.7);
\coordinate (v0) at (-2.100000,2.000000);
\coordinate (v1) at (0.000000,2.000000);
\coordinate (v2) at (2.100000,2.000000);
\coordinate (v3) at (-2.100000,-0.000000);
\coordinate (v4) at (-2.100000,-2.000000);
\coordinate (v5) at (0.000000,-0.000000);
\coordinate (v6) at (0.000000,-2.000000);
\coordinate (v7) at (2.100000,-0.000000);
\coordinate (v8) at (2.100000,-2.000000);
\draw[jfef edge] (v0) -- (v1);
\draw[jfef edge] (v0) .. controls (-0.700000,2.900000) and (0.700000,2.900000) .. (v2);
\draw[jfef edge] (v0) -- (v5);
\draw[jfef edge] (v0) -- (v6);
\draw[jfef edge] (v0) -- (v7);
\draw[jfef edge] (v0) .. controls (-0.079310,1.318391) and (1.320690,-0.014943) .. (v8);
\draw[jfef edge] (v1) -- (v2);
\draw[jfef edge] (v1) -- (v3);
\draw[jfef edge] (v1) -- (v4);
\draw[jfef edge] (v1) -- (v7);
\draw[jfef edge] (v1) -- (v8);
\draw[jfef edge] (v2) -- (v3);
\draw[jfef edge] (v2) .. controls (1.320690,0.014943) and (-0.079310,-1.318391) .. (v4);
\draw[jfef edge] (v2) -- (v5);
\draw[jfef edge] (v2) -- (v6);
\draw[jfef edge] (v3) -- (v4);
\draw[jfef edge] (v3) -- (v5);
\draw[jfef edge] (v3) -- (v6);
\draw[jfef edge] (v3) .. controls (-0.700000,0.900000) and (0.700000,0.900000) .. (v7);
\draw[jfef edge] (v3) -- (v8);
\draw[jfef edge] (v4) -- (v5);
\draw[jfef edge] (v4) -- (v6);
\draw[jfef edge] (v4) -- (v7);
\draw[jfef edge] (v4) .. controls (-0.700000,-1.100000) and (0.700000,-1.100000) .. (v8);
\draw[jfef edge] (v5) -- (v6);
\draw[jfef edge] (v5) -- (v7);
\draw[jfef edge] (v5) -- (v8);
\draw[jfef edge] (v6) -- (v7);
\draw[jfef edge] (v6) -- (v8);
\draw[jfef edge] (v7) -- (v8);
\node[jfef centre] at (v0) {$c_{0}$};
\node[jfef centre] at (v1) {$c_{1}$};
\node[jfef centre] at (v2) {$c_{2}$};
\node[jfef auxiliary] at (v3) {$x_{0,0}$};
\node[jfef auxiliary] at (v4) {$x_{1,0}$};
\node[jfef auxiliary] at (v5) {$x_{0,1}$};
\node[jfef auxiliary] at (v6) {$x_{1,1}$};
\node[jfef auxiliary] at (v7) {$x_{0,2}$};
\node[jfef auxiliary] at (v8) {$x_{1,2}$};
\end{tikzpicture}

\caption{The new vertices in three layers. Each column contains its centre
and two auxiliaries. The seed and its incident edges are omitted. A centre
is adjacent to every other centre and to the auxiliaries outside its own
column. All six auxiliaries are pairwise adjacent.}\label{fig:layers}
\end{figure}

\subsection{Deleting one vertex}\label{sec:consequences}
Can a single vertex deletion preserve the strict level of \(F_k\)?
The proof of strictness already provides two ways to examine the
remaining graph. After deleting a centre we have an explicit fold
reduction. For any other vertex, the Euler recurrence relates the value
of \(q\) on the remaining graph to the value on its neighbourhood graph.

\begin{proposition}\label{prop:single-deletion}
Let $k\geq2$ and let $F_k$ be the graph of
Theorem~\ref{thm:family3k8}. Deleting a centre gives a graph in $\D_0$.
Deleting any other vertex gives a graph outside every class $\D_j$,
$j\geq-1$. Consequently no single vertex deletion preserves its
strict level $k$.
\end{proposition}
\begin{proof}
Write \(F_k=G_m\) with \(m=k-1\geq1\). If we delete a centre
when \(m=1\), the remainder is \(B-p\), which folds to one vertex by
Lemma~\ref{lem:seed-rest}. For \(m\geq2\), Lemma~\ref{lem:remainder}
gives a fold reduction after deleting any centre. In both cases
the remainder belongs to \(\D_0\subseteq\D_{k-1}\) and thus cannot
have strict level \(k\).

Now let \(v\) be a vertex other than a centre. By
Lemma~\ref{lem:euler-links}, the non-zero values in
Table~\ref{tab:seed-cliques} remain non-zero as we add layers, so
\(q(F_k[N_{F_k}(v)])\ne0\). The graph \(F_k\) itself belongs to
\(\D_k\) by Theorem~\ref{thm:family3k8}, so
\(q(F_k)=0\) by Lemma~\ref{lem:qobstruction}. Using these two
values in the recurrence of Lemma~\ref{lem:qrec} gives
\[
 q(F_k-v)=q(F_k)+q(F_k[N_{F_k}(v)])
         =q(F_k[N_{F_k}(v)])\ne0.
\]
Lemma~\ref{lem:qobstruction} therefore excludes \(F_k-v\) from
every dismantling class.
\end{proof}

The conclusion concerns one deletion from \(F_k\), starting with
\(F_2=B\). Minimality among all proper induced subgraphs and least
order among all graphs of strict level \(k\) would require separate
arguments.


\runcheck{deletions}
The calculation checks every single-vertex deletion from $F_2,F_3,F_4$.
It verifies the fold reduction after a centre deletion and the non-zero
Euler value after any other deletion. Proposition~\ref{prop:single-deletion}
proves the dichotomy for every $k\geq2$.
